\subsection{সম্পর্ক ও কী: ধারণা যাচাই}
\textbf{নির্বাচিত MCQ}
\begin{chSixSideBySide}[.48]
\small 1. একটি বিভাগের একাধিক শিক্ষার্থী থাকলে বিভাগ থেকে শিক্ষার্থীর সম্পর্ক কোনটি?

\par
\begin{tabularx}{\linewidth}{@{}Y Y@{}}
ক. 1:1 & খ. 1:M \\
গ. M:N & ঘ. 0:1 \\
\end{tabularx}
\smallskip
\chSixMCQAnswer{খ}
\chSixNextColumn
\small 2. Many-to-Many Relationship বাস্তবায়নে সাধারণত কোনটি ব্যবহার করা হয়?

\par
\begin{tabularx}{\linewidth}{@{}Y Y@{}}
ক. Report & খ. Junction Table \\
গ. Index & ঘ. Form \\
\end{tabularx}
\smallskip
\chSixMCQAnswer{খ}
\end{chSixSideBySide}
\begin{chSixSideBySide}[.48]
\small 3. Child Table-এর যে ফিল্ড Parent-এর Key নির্দেশ করে তাকে কী বলে?

\par
\begin{tabularx}{\linewidth}{@{}Y Y@{}}
ক. Primary Key & খ. Foreign Key \\
গ. Sort Key & ঘ. Index Key \\
\end{tabularx}
\smallskip
\chSixMCQAnswer{খ}
\chSixNextColumn
\small 4. Parent Table-এ নেই এমন Foreign Key মান প্রতিরোধের সঙ্গে
কোন ধারণা সবচেয়ে বেশি সম্পর্কিত?

\par
\begin{tabularx}{\linewidth}{@{}Y Y@{}}
ক. Sorting & খ. Referential Integrity \\
গ. Report & ঘ. Form \\
\end{tabularx}
\smallskip
\chSixMCQAnswer{খ}
\end{chSixSideBySide}


\begin{practicebox}

\textbf{নিজে যাচাই করো}

একটি College ডেটাবেজে নিচের টেবিল আছে---

\noindent\texttt{Department(DeptID, DeptName)}\par

\noindent\texttt{Student(StudentID, Name, DeptID)}\par

\noindent\texttt{Result(ResultID, StudentID, ICTMark)}\par

প্রতিটি শিক্ষার্থী একটি Department-এর অন্তর্ভুক্ত। একজন শিক্ষার্থীর
একাধিক পরীক্ষার ফল আলাদা ResultID দিয়ে সংরক্ষিত হতে পারে।

প্রশ্ন:

\begin{enumerate}
    \item কোন ফিল্ডগুলো Primary Key হতে পারে?
    \item কোন ফিল্ডগুলো Foreign Key?
    \item Department ও Student-এর সম্পর্ক কী?
    \item Student ও Result-এর সম্পর্ক কী?
    \item Result থেকে Student-এর সম্পর্ক কোন দিক থেকে কী নামে পরিচিত?
\end{enumerate}

\end{practicebox}





\par\nointerlineskip\vspace{4pt}\noindent\begin{minipage}{\linewidth}
\subsection{সমন্বিত CQ: অনুমতি ও পুনরুদ্ধার}
কলেজের ফলাফলের ডেটাবেজের সর্বশেষ ব্যাকআপ নেওয়া হয়েছিল সোমবার বিকেল ৫টায়। মঙ্গলবার নতুন ফল যোগ করার পর ডিস্ক নষ্ট হয়ে যায়। শিক্ষক নিজের বিষয়ের ফল সংশোধন করতে পারেন; প্রকাশনা কর্মকর্তা শুধু অনুমোদিত ফল পড়তে পারেন। সোমবারের ব্যাকআপ ব্যবহার করে ডেটাবেজটি পুনরুদ্ধার করা হলো।
\textbf{প্রশ্ন}
\begin{enumerate}[label=(\alph*)]
\item ব্যাকআপ কী?
\item সঠিক পাসওয়ার্ড থাকলেও প্রকাশনা কর্মকর্তা ফল বদলাতে পারেন না কেন?
\item ডেটাবেজ ফিরিয়ে আনার পরে কোন তথ্যগুলো যাচাই করতে হবে?
\item সোমবারের ব্যাকআপ-ই মঙ্গলবারের সব তথ্য ফেরানোর নিশ্চয়তা দেয়—বিশ্লেষণ করো।
\end{enumerate}
\end{minipage}\par
% Preserved practice moved from theory to the end-of-chapter bank.

\subsection{Data Security: ধারণা যাচাই}
\begin{multicols}{2}
\small
\begin{enumerate}
\item \begin{minipage}[t]{\linewidth}
একজন শিক্ষার্থী নিজের ফল দেখতে পারে, কিন্তু পরিবর্তন
          করতে পারে না। তাকে কোন ধরনের Authorization দেওয়া হয়েছে?
\end{minipage}
\item \begin{minipage}[t]{\linewidth}
একজন শিক্ষক সঠিক পরিচয়ে লগইন করেছেন, কিন্তু অন্য বিভাগের ফল
          পরিবর্তনের অনুমতি পাননি। এই কাজের সীমা নির্ধারণে Authentication না
          Authorization কাজ করছে?
\end{minipage}
\item \begin{minipage}[t]{\linewidth}
ডিস্ক নষ্ট হয়ে ফলাফলের ডেটাবেজ হারিয়ে গেছে।
          পাসওয়ার্ড জানা থাকলে কি হারানো ডেটা ফিরে পাওয়া যাবে? কেন?
\end{minipage}
\item \begin{minipage}[t]{\linewidth}
ডেটাবেজে GPA-এর রেকর্ড আছে, কিন্তু সংরক্ষিত মান ভুল হয়ে গেছে।
          এটি Data Loss নাকি Data Corruption?
\end{minipage}
\item \begin{minipage}[t]{\linewidth}
কে কখন ফল পরিবর্তন করেছে তা জানতে নিরাপত্তার কোন ব্যবস্থা
          সবচেয়ে প্রাসঙ্গিক?
\end{minipage}
\item \begin{minipage}[t]{\linewidth}
Mobile Number
          \texttt{017*****321}
          আকারে দেখানো হলে কোন নিরাপত্তাপদ্ধতি ব্যবহার করা হয়েছে?
\end{minipage}
\end{enumerate}
\end{multicols}



\textbf{নির্বাচিত MCQ}
\begin{chSixSideBySide}[.48]
\small 1. ব্যবহারকারীর পরিচয় যাচাই করে কোনটি?

\par
\begin{tabularx}{\linewidth}{@{}Y Y@{}}
ক. Authorization & খ. Authentication \\
গ. Recovery & ঘ. Masking \\
\end{tabularx}
\smallskip
\chSixMCQAnswer{খ}
\chSixNextColumn
\small 2. ডেটা শুধু পড়ার অনুমতি কোনটি?

\par
\begin{tabularx}{\linewidth}{@{}Y Y@{}}
ক. Read Authorization & খ. Drop Authorization \\
গ. Alteration Authorization & ঘ. Resource Authorization \\
\end{tabularx}
\smallskip
\chSixMCQAnswer{ক}
\end{chSixSideBySide}
\begin{chSixSideBySide}[.48]
\small 3. সংরক্ষণের ডিস্ক নষ্ট হয়ে ডেটা হারানো কোন ধরনের ব্যর্থতার উদাহরণ?

\par
\begin{tabularx}{\linewidth}{@{}Y Y@{}}
ক. Transaction Failure & খ. Disk Failure \\
গ. Authorization Failure & ঘ. Read Failure \\
\end{tabularx}
\smallskip
\chSixMCQAnswer{খ}
\chSixNextColumn
\small 4. কে কখন রেকর্ড পরিবর্তন করেছে তা জানতে সবচেয়ে উপযোগী কোনটি?

\par
\begin{tabularx}{\linewidth}{@{}Y Y@{}}
ক. Data Masking & খ. Audit Log \\
গ. Drop Authorization & ঘ. Sorting \\
\end{tabularx}
\smallskip
\chSixMCQAnswer{খ}
\end{chSixSideBySide}


\Needspace{7\baselineskip}

\subsection{ডেটা এনক্রিপশন: ধারণা যাচাই}
\begin{multicols}{2}
\small
\begin{enumerate}
\item \begin{minipage}[t]{\linewidth}
Plaintext এবং Ciphertext-এর পার্থক্য কী?
\end{minipage}
\item \begin{minipage}[t]{\linewidth}
Symmetric এনক্রিপশনে একই কী ব্যবহারের প্রধান
          সুবিধা ও সমস্যা কী?
\end{minipage}
\item \begin{minipage}[t]{\linewidth}
Asymmetric পদ্ধতিতে Public Key ও Private Key-এর কাজ কী?
\end{minipage}
\item \begin{minipage}[t]{\linewidth}
Caesar Cipher-এ \(k=3\) হলে \texttt{ABC} কী হবে?
\end{minipage}
\item \begin{minipage}[t]{\linewidth}
এনক্রিপ্ট করা ব্যাকআপ ফাইলের কী হারিয়ে গেলে কী সমস্যা হবে?
\end{minipage}
\item \begin{minipage}[t]{\linewidth}
কেন এনক্রিপশনকে ব্যাকআপের বিকল্প বলা যায় না?
\end{minipage}
\end{enumerate}
\end{multicols}



\textbf{নির্বাচিত MCQ}
\begin{chSixSideBySide}[.48]
\small 1. এনক্রিপশনের আগে মূল বার্তাকে কী বলা হয়?

\par
\begin{tabularx}{\linewidth}{@{}Y Y@{}}
ক. Ciphertext & খ. Plaintext \\
গ. Key & ঘ. Algorithm \\
\end{tabularx}
\smallskip
\chSixMCQAnswer{খ}
\chSixNextColumn
\small 2. একই Secret Key দিয়ে এনক্রিপশন ও ডিক্রিপশন কোন পদ্ধতিতে হয়?

\par
\begin{tabularx}{\linewidth}{@{}Y Y@{}}
ক. Asymmetric & খ. Symmetric \\
গ. Public-only & ঘ. Hashing \\
\end{tabularx}
\smallskip
\chSixMCQAnswer{খ}
\end{chSixSideBySide}
\begin{chSixSideBySide}[.48]
\small 3. কোনটি Asymmetric Algorithm-এর পরিচিত উদাহরণ?

\par
\begin{tabularx}{\linewidth}{@{}Y Y@{}}
ক. AES & খ. DES \\
গ. RSA & ঘ. IDEA \\
\end{tabularx}
\smallskip
\chSixMCQAnswer{গ}
\chSixNextColumn
\small 4. কোন Key গোপন রাখতে হয়?

\par
\begin{tabularx}{\linewidth}{@{}Y Y@{}}
ক. Public Key & খ. Private Key \\
গ. দুটিই প্রকাশ্য & ঘ. কোনোটিই নয় \\
\end{tabularx}
\smallskip
\chSixMCQAnswer{খ}
\end{chSixSideBySide}
